% SPDX-License-Identifier: Apache-2.0
% covers: SetTable
\datasheet{SetTable}{A set-associative key-to-value table in block RAM:
lookup, insert, update and remove, one operation a cycle.}

\dsfacts{\code{lib/parts/src/settable.rs}}{\code{txhdl\_parts::settable::SetTable}}%
{\code{//docs:datasheets}}

\subsection*{Function}

\code{SetTable} keeps keys and their values in \code{WAYS} ways of
\code{ROWS} rows each. A key's set, the row it may live in, is chosen by
a multiplicative hash, and the key may sit in any way of that set.

It takes operations on \code{Rx<TableOp<KEY, VALUE>>}: \code{op}, one of
\code{LOOKUP}, \code{INSERT}, \code{UPDATE} and \code{REMOVE}; the
\code{key}; and the \code{val} an insert or an update puts in. It
answers each on \code{Tx<TableResult<VALUE>>}: \code{res}, one of
\code{HIT}, \code{MISS}, \code{EXISTS} and \code{FULL}, and a
\code{val}. Its other output, \code{Out<U<32>>}, counts the inserts
refused because their set was full.

\code{SetTableModel} is the same table in plain Rust, for a design's
own tests: the same hash, the same way chosen and the same refusals.
\code{set\_of} is the one definition of the hash, which the part and
the model both use.

\subsection*{Parameters}

\begin{center}\footnotesize
\begin{tabular}{@{}lp{0.7\textwidth}@{}}
\toprule
Parameter & Meaning \\
\midrule
\code{KEY} & The key's width, at most 32 bits. \\
\code{VALUE} & The value's width. \\
\code{ROW\_ADDR} & The bits that number the sets, from 1 to 32. \\
\code{ROWS} & The rows of a way, which must be \code{1 << ROW\_ADDR}. \\
\code{WAYS} & The ways, at most 256. \\
\bottomrule
\end{tabular}
\end{center}

Stable Rust cannot write \code{ROWS} from \code{ROW\_ADDR} in a type, so both are
given; the plain-Rust \code{set\_of} and \code{SetTableModel::new}
check \code{ROW\_ADDR}. The tables below use \code{SetTable<16, 8, 3, 8, 4>},
the table of the tests: eight sets of four ways, keys of 16 bits and
values of 8.

\dstables{SetTable}

\subsection*{Behaviour}

\begin{itemize}
\item \textbf{The set.} A key's set is \code{(key * 0x9E3779B1) >> (32
- B)} over the key widened to 32 bits: the top \code{ROW\_ADDR} bits of the
product, a Fibonacci hash. A key's low bits would crowd keys allocated
in strides into a few sets, and the product spreads them.
\item \textbf{The ways.} Each way is three memories: its keys, its
values and whether each row is used, each a \code{Mems} of \code{WAYS}.
An operation reads every way at its set at once, into registers at the
edge, the synchronous read a block RAM has. The cycle after, a loop
over the ways finds a hit, the first free way and whether the set is
full, and answers; the one way that changes is written back in the
cycle the answer goes out, which is that cycle unless the answer
waits. So every
memory is read at one address and written at one, which a block RAM
holds.
\item \textbf{One operation a cycle.} An operation that reads the set
the one before it is writing, in the same cycle, would read the row
from before that write. So the way written is forwarded into the row
it reads, and nothing stalls.
\item \textbf{Lookup} answers the value of a key present with
\code{HIT}, or \code{MISS}.
\item \textbf{Insert} puts a key not present into the lowest-numbered
free way of its set and answers \code{HIT}. A key present changes
nothing and answers \code{EXISTS} with the value it already has. A
full set takes nothing, answers \code{FULL}, and counts one refusal.
\item \textbf{Update} replaces the value of a key present and answers
the new value with \code{HIT}, or \code{MISS}. A function of the old
value is the user's: look up, work it out, update.
\item \textbf{Remove} frees a key present and answers its value with
\code{HIT}, or \code{MISS}.
\item An answer's \code{val} is, on \code{HIT}, the value found, put in
or freed; on \code{EXISTS}, the value the key already has; and on
\code{MISS} or \code{FULL}, nought.
\item An operation is taken when its answer has room: when nothing is
being answered, or the answer is taken that cycle. Its write is made
once, in the cycle its answer goes out, however long the answer waits.
\end{itemize}

\subsection*{Verification}

\begin{itemize}
\item \code{//lib/parts:parts\_test} runs \code{settable::tests}:
\begin{itemize}
\item \code{random\_operations\_agree\_with\_the\_model}: 2000
operations from a fixed seed on 48 keys, so that every answer occurs
more than twenty times, against \code{SetTableModel}. The operations
are offered on three cycles in four and the answers held up for up to
eight cycles at a time, and the refusals counted must agree too.
\item \code{back\_to\_back\_on\_one\_set\_see\_each\_other}: six keys of
one set, each inserted and looked up back to back, then a remove, an
insert, an update and a lookup, each reading the row the one before it
is writing; two inserts find the set full.
\item \code{strided\_keys\_spread\_over\_the\_sets}: 256 keys at strides
of 64 and of 4096 land in at least 48 of 64 sets.
\item \code{each\_way\_is\_a\_memory\_read\_through\_a\_register} reads the
Verilog of \code{SetTable<16, 8, 3, 8, 4>}.
\item \code{a\_large\_table\_fits\_block\_rams} lowers
\code{SetTable<32, 32, 10, 1024, 4>}, whose values are 1024 rows of 32
bits a way, through the check that a memory of more than 4096 bits
fits a block RAM's ports (issue 1285).
\item \code{the\_parameters\_agree}: the tests' table has \code{ROWS} as
\code{1 << ROW\_ADDR}, and \code{set\_of} puts each of the first 1000 keys in
one of its eight sets.
\end{itemize}
\item Its netlist is not simulated against a trace.
\end{itemize}

\subsection*{Used by}

Nothing in the tree yet, outside its own tests.

\subsection*{Limits}

A key is at most 32 bits, since the hash is over 32. \code{WAYS} is at
most 256, since a way's number is eight bits. \code{ROWS} must be
\code{1 << ROW\_ADDR}, given beside it. A set holds at most \code{WAYS} keys, and
an insert beyond that is refused rather than placed elsewhere.
